\documentclass[10pt,a4paper]{amsart}
\usepackage{amsmath,amssymb,amsthm,graphicx}
\usepackage{verbatim}
\newenvironment{reference}{\comment}{\endcomment}
%\newenvironment{reference}{}{}
\newenvironment{slogan}{\comment}{\endcomment}
\newenvironment{history}{\comment}{\endcomment}

% For commutative diagrams we use Xy-pic
\usepackage[all]{xy}

% We use 2cell for 2-commutative diagrams.
\xyoption{2cell}
% Theorem environments.
%
\theoremstyle{plain}
\newtheorem{theorem}[subsection]{Theorem}
\newtheorem{proposition}[subsection]{Proposition}
\newtheorem{lemma}[subsection]{Lemma}

\theoremstyle{definition}
\newtheorem{definition}[subsection]{Definition}
\newtheorem{example}[subsection]{Example}
\newtheorem{exercise}[subsection]{Exercise}
\newtheorem{situation}[subsection]{Situation}

\theoremstyle{remark}
\newtheorem{remark}[subsection]{Remark}
\newtheorem{remarks}[subsection]{Remarks}

\numberwithin{equation}{subsection}

% Macros
%
\def\lim{\mathop{\mathrm{lim}}\nolimits}
\def\colim{\mathop{\mathrm{colim}}\nolimits}
\def\Spec{\mathop{\mathrm{Spec}}}
\def\Hom{\mathop{\mathrm{Hom}}\nolimits}
\def\Ext{\mathop{\mathrm{Ext}}\nolimits}
\def\SheafHom{\mathop{\mathcal{H}\!\mathit{om}}\nolimits}
\def\SheafExt{\mathop{\mathcal{E}\!\mathit{xt}}\nolimits}
\def\Sch{\mathit{Sch}}
\def\Mor{\mathop{\mathrm{Mor}}\nolimits}
\def\Ob{\mathop{\mathrm{Ob}}\nolimits}
\def\Sh{\mathop{\mathit{Sh}}\nolimits}
\def\NL{\mathop{N\!L}\nolimits}
\def\CH{\mathop{\mathrm{CH}}\nolimits}
\def\proetale{{pro\text{-}\acute{e}tale}}
\def\etale{{\acute{e}tale}}
\def\QCoh{\mathit{QCoh}}
\def\Ker{\mathop{\mathrm{Ker}}}
\def\Im{\mathop{\mathrm{Im}}}
\def\Coker{\mathop{\mathrm{Coker}}}
\def\Coim{\mathop{\mathrm{Coim}}}

% Boxtimes
%
\DeclareMathSymbol{\boxtimes}{\mathbin}{AMSa}{"02}

%
% Macros for moduli stacks/spaces
%
\def\QCohstack{\mathcal{QC}\!\mathit{oh}}
\def\Cohstack{\mathcal{C}\!\mathit{oh}}
\def\Spacesstack{\mathcal{S}\!\mathit{paces}}
\def\Quotfunctor{\mathrm{Quot}}
\def\Hilbfunctor{\mathrm{Hilb}}
\def\Curvesstack{\mathcal{C}\!\mathit{urves}}
\def\Polarizedstack{\mathcal{P}\!\mathit{olarized}}
\def\Complexesstack{\mathcal{C}\!\mathit{omplexes}}
% \Pic is the operator that assigns to X its picard group, usage \Pic(X)
% \Picardstack_{X/B} denotes the Picard stack of X over B
% \Picardfunctor_{X/B} denotes the Picard functor of X over B
\def\Pic{\mathop{\mathrm{Pic}}\nolimits}
\def\Picardstack{\mathcal{P}\!\mathit{ic}}
\def\Picardfunctor{\mathrm{Pic}}
\def\Deformationcategory{\mathcal{D}\!\mathit{ef}}
\usepackage{cleveref}
\crefname{originallemma}{original lemma}{original lemmas}
\Crefname{originallemma}{Original lemma}{Original lemmas}
\makeatletter
\newlabel{lemma-twisted-inverse-image}{{0A9E}{source}{}{sourcepointer-1}{}}
\newlabel{lemma-twisted-inverse-image@cref}{{[originallemma][0][]0A9E}{[1][0][]source}{}{}{}}
\makeatother
\crefname{originalreference}{original reference}{original references}
\Crefname{originalreference}{Original reference}{Original references}
\makeatletter
\newlabel{remark-relative-dualizing-complex}{{0B6S}{source}{}{sourcepointer-2}{}}
\newlabel{remark-relative-dualizing-complex@cref}{{[originalreference][0][]0B6S}{[1][0][]source}{}{}{}}
\makeatother
\crefname{originallemma}{original lemma}{original lemmas}
\Crefname{originallemma}{Original lemma}{Original lemmas}
\makeatletter
\newlabel{divisors-lemma-immersion-smooth-into-smooth-regular-immersion}{{067U}{source}{}{sourcepointer-3}{}}
\newlabel{divisors-lemma-immersion-smooth-into-smooth-regular-immersion@cref}{{[originallemma][0][]067U}{[1][0][]source}{}{}{}}
\makeatother
\crefname{originallemma}{original lemma}{original lemmas}
\Crefname{originallemma}{Original lemma}{Original lemmas}
\makeatletter
\newlabel{more-morphisms-lemma-regular-immersion-perfect}{{068C}{source}{}{sourcepointer-4}{}}
\newlabel{more-morphisms-lemma-regular-immersion-perfect@cref}{{[originallemma][0][]068C}{[1][0][]source}{}{}{}}
\makeatother
\crefname{originalequation}{original equation}{original equations}
\Crefname{originalequation}{Original equation}{Original equations}
\makeatletter
\newlabel{equation-pre-rigid}{{0E2Q}{source}{}{sourcepointer-5}{}}
\newlabel{equation-pre-rigid@cref}{{[originalequation][0][]0E2Q}{[1][0][]source}{}{}{}}
\makeatother
\crefname{originallemma}{original lemma}{original lemmas}
\Crefname{originallemma}{Original lemma}{Original lemmas}
\makeatletter
\newlabel{lemma-compare-with-pullback-flat-proper-noetherian}{{0A9U}{source}{}{sourcepointer-6}{}}
\newlabel{lemma-compare-with-pullback-flat-proper-noetherian@cref}{{[originallemma][0][]0A9U}{[1][0][]source}{}{}{}}
\makeatother
\crefname{originallemma}{original lemma}{original lemmas}
\Crefname{originallemma}{Original lemma}{Original lemmas}
\makeatletter
\newlabel{lemma-regular-immersion}{{0BR0}{source}{}{sourcepointer-7}{}}
\newlabel{lemma-regular-immersion@cref}{{[originallemma][0][]0BR0}{[1][0][]source}{}{}{}}
\makeatother
\crefname{originallemma}{original lemma}{original lemmas}
\Crefname{originallemma}{Original lemma}{Original lemmas}
\makeatletter
\newlabel{morphisms-lemma-differentials-diagonal}{{08S2}{source}{}{sourcepointer-8}{}}
\newlabel{morphisms-lemma-differentials-diagonal@cref}{{[originallemma][0][]08S2}{[1][0][]source}{}{}{}}
\makeatother
\title{For a smooth proper morphism the right adjoint of pushforward is the shift of the top exterior power of the sheaf of differentials}
\author{The Stacks Project Authors}
\date{stacks-project:duality.tex:local-snapshot-20261012; GNU FDL 1.2 or later}
\begin{document}
\maketitle
\noindent\textit{Editorial note.} All delivered mathematical text is the author's own, copied from the pinned Stacks Project chapter file duality.tex; no statement or proof step was rewritten, reordered or removed.\par\medskip
\noindent\textit{Editorial note.} Preamble provenance: the pinned chapter file opens with an input of the shared preamble and carries no preamble text of its own, so the project's single Stacks-wide preamble staged at sources/stacks-duality/preamble.tex is excerpted verbatim under the recipe fixed by LANE\_BRIEF 22.1, as three spans, lines 8-13 (verbatim package and the reference/slogan/history environments), lines 14-18 (Xy-pic setup) and lines 154-217 (the native theorem declarations and mathematical macros). Nothing was added, removed or redefined; the wrapper class and the amsmath/amssymb/amsthm/graphicx base packages are supplied by the assembler.\par\medskip
\noindent\textit{Editorial note.} The counted statement and proof contain no cross reference, no citation and no included file, so the delivered body is self contained.\par\medskip
\noindent\textit{Editorial note.} Canonical location of the counted proposition: Stacks tag 0BRT (sources/stacks-duality/pinned-original-source.tex line 3879).\par\medskip
\noindent\textit{Allowed standard prerequisites (editorial).}\begin{itemize}
\item The right adjoint of pushforward and the pre-rigid identity. Same chapter, Lemma labelled lemma-twisted-inverse-image, invoked in the statement, and equation labelled equation-pre-rigid, cited at proof lines 3902-3906. Fixes a and provides the first equality.
\item The diagonal of a smooth morphism is a Koszul-regular immersion and hence perfect proper. Divisors, Lemma labelled divisors-lemma-immersion-smooth-into-smooth-regular-immersion, cited at proof lines 3896-3898, and More on Morphisms, Lemma labelled more-morphisms-lemma-regular-immersion-perfect, cited at proof lines 3898-3900. Makes c available and the computation meaningful.
\item The comparison with pullback for a flat proper Noetherian morphism and the regular immersion computation. Same chapter, Lemmas labelled lemma-compare-with-pullback-flat-proper-noetherian and lemma-regular-immersion, cited at proof lines 3908-3918. Give the second and fourth equalities of the chain.
\item The conormal sheaf of the diagonal is the sheaf of differentials. Morphisms, Lemma labelled morphisms-lemma-differentials-diagonal, cited at proof lines 3928-3930. Identifies C\_Delta with Omega\_\{X/S\}, finishing the proof.
\end{itemize}\medskip
\ifdefined\hypertarget\hypertarget{sourcepointer-1}{}\fi
\noindent\textit{Source reference.} Original lemma 0A9E: sources/stacks-duality/pinned-original-source.tex, source lines 653--653. Same-chapter pointer into the author environment lemma at source line 653; invoked at proof line 3894 of the counted proof. It is retained only as a navigation pointer; every load-bearing step is inside the delivered span.\par\medskip
\ifdefined\hypertarget\hypertarget{sourcepointer-2}{}\fi
\noindent\textit{Source reference.} Original reference 0B6S: sources/stacks-duality/pinned-original-source.tex, source lines 2945--2945. Same-chapter pointer into the author environment remark at source line 2945; invoked at proof line 3893 of the counted proof. It is retained only as a navigation pointer; every load-bearing step is inside the delivered span.\par\medskip
\ifdefined\hypertarget\hypertarget{sourcepointer-3}{}\fi
\noindent\textit{Source reference.} Original lemma 067U: sources/stacks-chow/tags.txt, source lines 7983--7983. Cross-chapter standard prerequisite located in the project tag list staged in the local snapshot; invoked at proof line 3898 of the counted proof. It is retained only as a navigation pointer; every load-bearing step is inside the delivered span.\par\medskip
\ifdefined\hypertarget\hypertarget{sourcepointer-4}{}\fi
\noindent\textit{Source reference.} Original lemma 068C: sources/stacks-chow/tags.txt, source lines 8001--8001. Cross-chapter standard prerequisite located in the project tag list staged in the local snapshot; invoked at proof line 3900 of the counted proof. It is retained only as a navigation pointer; every load-bearing step is inside the delivered span.\par\medskip
\ifdefined\hypertarget\hypertarget{sourcepointer-5}{}\fi
\noindent\textit{Source reference.} Original equation 0E2Q: sources/stacks-duality/pinned-original-source.tex, source lines 3166--3166. Same-chapter pointer into the author environment equation at source line 3166; invoked at proof line 3918 of the counted proof. It is retained only as a navigation pointer; every load-bearing step is inside the delivered span.\par\medskip
\ifdefined\hypertarget\hypertarget{sourcepointer-6}{}\fi
\noindent\textit{Source reference.} Original lemma 0A9U: sources/stacks-duality/pinned-original-source.tex, source lines 3306--3306. Same-chapter pointer into the author environment lemma at source line 3306; invoked at proof line 3920 of the counted proof. It is retained only as a navigation pointer; every load-bearing step is inside the delivered span.\par\medskip
\ifdefined\hypertarget\hypertarget{sourcepointer-7}{}\fi
\noindent\textit{Source reference.} Original lemma 0BR0: sources/stacks-duality/pinned-original-source.tex, source lines 3852--3852. Same-chapter pointer into the author environment lemma at source line 3852; invoked at proof line 3922 of the counted proof. It is retained only as a navigation pointer; every load-bearing step is inside the delivered span.\par\medskip
\ifdefined\hypertarget\hypertarget{sourcepointer-8}{}\fi
\noindent\textit{Source reference.} Original lemma 08S2: sources/stacks-chow/tags.txt, source lines 11115--11115. Cross-chapter standard prerequisite located in the project tag list staged in the local snapshot; invoked at proof line 3929 of the counted proof. It is retained only as a navigation pointer; every load-bearing step is inside the delivered span.\par\medskip
\begin{lemma}
\label{lemma-smooth-proper}
Let $S$ be a Noetherian scheme.
Let $f : X \to S$ be a smooth proper morphism of relative dimension $d$.
Let $a$ be the right adjoint of
$Rf_* : D_\QCoh(\mathcal{O}_X) \to D_\QCoh(\mathcal{O}_S)$ as in
Lemma \ref{lemma-twisted-inverse-image}. Then there is an isomorphism
$$
\wedge^d \Omega_{X/S}[d] \longrightarrow a(\mathcal{O}_S)
$$
in $D(\mathcal{O}_X)$.
\end{lemma}

\begin{proof}
Set $\omega_{X/S}^\bullet = a(\mathcal{O}_S)$ as in
Remark \ref{remark-relative-dualizing-complex}.
Let $c$ be the right adjoint of Lemma \ref{lemma-twisted-inverse-image} for
$\Delta : X \to X \times_S X$. Because $\Delta$
is the diagonal of a smooth morphism it is a
Koszul-regular immersion, see Divisors, Lemma
\ref{divisors-lemma-immersion-smooth-into-smooth-regular-immersion}.
In particular, $\Delta$ is a perfect proper morphism
(More on Morphisms, Lemma \ref{more-morphisms-lemma-regular-immersion-perfect})
and we obtain
\begin{align*}
\mathcal{O}_X
& =
c(L\text{pr}_1^*\omega_{X/S}^\bullet) \\
& =
L\Delta^*(L\text{pr}_1^*\omega_{X/S}^\bullet)
\otimes_{\mathcal{O}_X}^\mathbf{L}
c(\mathcal{O}_{X \times_S X}) \\
& =
\omega_{X/S}^\bullet \otimes_{\mathcal{O}_X}^\mathbf{L}
c(\mathcal{O}_{X \times_S X}) \\
& =
\omega_{X/S}^\bullet
\otimes_{\mathcal{O}_X}^\mathbf{L}
\wedge^d(\mathcal{N}_\Delta)[-d]
\end{align*}
The first equality is (\ref{equation-pre-rigid}) because
$\omega_{X/S}^\bullet = a(\mathcal{O}_S)$. The second equality by
Lemma \ref{lemma-compare-with-pullback-flat-proper-noetherian}.
The third equality because $\text{pr}_1 \circ \Delta = \text{id}_X$.
The fourth equality by Lemma \ref{lemma-regular-immersion}.
Observe that $\wedge^d(\mathcal{N}_\Delta)$ is an invertible
$\mathcal{O}_X$-module. Hence $\wedge^d(\mathcal{N}_\Delta)[-d]$
is an invertible object of $D(\mathcal{O}_X)$ and we conclude that
$a(\mathcal{O}_S) = \omega_{X/S}^\bullet = \wedge^d(\mathcal{C}_\Delta)[d]$.
Since the conormal sheaf $\mathcal{C}_\Delta$ of $\Delta$ is
$\Omega_{X/S}$ by
Morphisms, Lemma \ref{morphisms-lemma-differentials-diagonal}
the proof is complete.
\end{proof}

\end{document}
